\section{八屏障与结构极小元}\label{sec:eight-barrier}

现在证明定理~\ref{thm:eight-barrier-trichotomy}，不依赖 17 个逐一列出的竞争者
证书。证明从平方算子的一个精确局部恒等式出发，最后归结为两缺陷几何的
有限四情形分析。

\subsection{算子平方}

从式~\eqref{eq:periodic-operator} 出发，把四种跃迁相乘两次。从剩余类 \(i\)
出发，\(A_\tau^2\) 的完整一行为

\begin{table}[tbp]
\centering
\caption{平方算子按位移分类的系数。}
\label{tab:squared-operator-coefficients}
\begingroup
\renewcommand{\arraystretch}{1.12}
\begin{tabular}{ll}
\toprule
位移 & 系数 \\
\midrule
\(-4\) & \(Q_{i-4}Q_{i-3}\) \\
\(-3\) & \(\tau_{i-3}(1+Q_{i-3})\) \\
\(-2\) & \(1\) \\
\(-1\) & \(\tau_{i-2}(1+Q_{i-2})\) \\
\(0\) & \(4\) \\
\(+1\) & \(\tau_{i-1}(1+Q_{i-1})\) \\
\(+2\) & \(1\) \\
\(+3\) & \(\tau_i(1+Q_i)\) \\
\(+4\) & \(Q_iQ_{i+1}\) \\
\bottomrule
\end{tabular}
\endgroup
\end{table}

为完整起见，以位移 \(+3\) 为例。它来自路径 \(+1,+2\) 和 \(+2,+1\)，
权重分别为 \(\tau_{i+1}\) 与 \(\tau_i\)。由于
\(\tau_{i+1}=\tau_iQ_i\)，二者之和为 \(\tau_i(1+Q_i)\)。其他奇位移由
平移和反向得到。位移 \(+4\) 只有路径 \(+2,+2\)，权重为
\(\tau_i\tau_{i+2}=Q_iQ_{i+1}\)。对角线上有四次立即折返；位移
\(\pm2\) 在其余路径消去后留下单位系数；反向给出负位移。

所以负通量 \(Q_i=-1\) 消去相关的奇位移耦合，而正通量以绝对值为二的振幅
打开这些耦合。这就是定理背后的缺陷机制。

\subsection{全负基线}\label{subsec:all-negative-baseline}

若对每个 \(i\) 都有 \(Q_i=-1\)，选择提升 \(\tau_i=(-1)^i\)。令 \(S\) 为
双边平移，并置

\begin{equation}
C=S+S^{-1},       E=S^{2}+S^{-2},       D=\operatorname{diag}((-1)^i).
\end{equation}

则 \(A=C+DE\)，同时 \(CD=-DC\)、\(ED=DE\)。因此交叉项消去：

\begin{equation}
\label{eq:all-negative-square}
A^{2}=C^{2}+E^{2}=4I+S^{2}+S^{-2}+S^{4}+S^{-4}.
\end{equation}

在 Fourier 模 \(S=e^{i\theta}\) 上，式~\eqref{eq:all-negative-square} 的符号为

\begin{equation}
\label{eq:all-negative-dispersion}
4+2\cos(2\theta)+2\cos(4\theta)\leq 8.
\end{equation}

等号在 \(\theta=0\) 处成立。因此全负相位满足 \(R=8\)。

\subsection{周期八的前三个矩}

对周期八字，式~\eqref{eq:moment-and-excess}--\eqref{eq:moment-integral} 给出带
符号闭步矩。\(A^2\) 在八个剩余类上的对角元均为四，故

\begin{equation}
\label{eq:period-eight-first-moment}
M_{1}=32.
\end{equation}

由于 \(A^2\) 是 Hermitian 矩阵，

\begin{equation}
M_{2}=\operatorname{tr}(A^{4})=\sum_{i,j}\lvert (A^{2})_{i,j}\rvert^{2}.
\end{equation}

对角项和偶位移项贡献 160。每个正通量打开四个定向奇位移矩阵元，每个元素
的平方为四，所以

\begin{equation}
\label{eq:period-eight-second-moment}
M_{2}=160+16d.
\end{equation}

计算 \(M_3=\operatorname{tr}((A^2)^3)\) 时，按局部闭乘积不使用激活点、使用
一个激活点、使用相邻激活点对或使用距离二激活点对分类。直接乘上表所示的
\(A^2\) 行，分别得到系数 \(944,168,96,48\)：

\begin{equation}
\label{eq:period-eight-third-moment}
M_{3}=944+168d+96a+48b.
\end{equation}

第~\ref{sec:general-period} 节通过枚举长度六的 430 个闭步字，给出式~
\eqref{eq:period-eight-third-moment} 的独立推导，且该推导对每个周期有效。因此
式~\eqref{eq:period-eight-first-moment}--\eqref{eq:period-eight-third-moment}
不涉及任何谱近似。

第二超额量为

\begin{equation}
\label{eq:period-eight-second-excess}
F_{2}=M_{3}-8M_{2}=-336+40d+96a+48b.
\end{equation}

\subsection{四个或更多缺陷}\label{subsec:four-or-more-defects}

由于 \(\prod_iQ_i=1\)，缺陷数 \(d\) 为偶数。

设 \(d=4\)。把四个正号之间的循环间隔写成和为八的正整数。若至少两个间隔
为一，则 \(2a\geq4\)。若没有间隔为一，则四个间隔全为二，且 \(b=4\)。
若恰有一个间隔为一，则其他间隔为 \(2,2,3\)，仍有 \(2a+b=4\)。所以总有
\(2a+b\geq4\)，式~\eqref{eq:period-eight-second-excess} 给出

\begin{equation}
F_{2}=-176+48(2a+b)\geq 16>0.
\end{equation}

若 \(d=6\)，两个负号位置至多破坏八条相邻正号边中的四条，故 \(a\geq4\)，
于是 \(F_2\geq288>0\)。若 \(d=8\)，式~
\eqref{eq:period-eight-first-moment}--\eqref{eq:period-eight-second-moment} 给出

\begin{equation}
F_{1}=M_{2}-8M_{1}=32>0.
\end{equation}

由引理~\ref{lem:moment-barrier}，每个至少有四个缺陷的相位都满足 \(R(Q)>8\)。

\subsection{两个缺陷}\label{subsec:two-defects}

令 \(d=2\)，以 \(s\in\{1,2,3,4\}\) 表示两个缺陷的较小循环间距。平移和
反射使 \(s\) 成为完整轨道不变量。精确闭步枚举给出下列首个正超额量：

\begin{table}[tbp]
\centering
\caption{两个非对径缺陷的首个正矩超额量。}
\label{tab:two-defect-excesses}
\begingroup
\renewcommand{\arraystretch}{1.12}
\begin{tabular}{lll}
\toprule
间距 \(s\) & 首个正超额量 & 数值 \\
\midrule
1 & \(F_{4}\) & 5,504 \\
2 & \(F_{6}\) & 64,336 \\
3 & \(F_{9}\) & 2,872,096 \\
\bottomrule
\end{tabular}
\endgroup
\end{table}

下面说明这些数如何认证。从八个剩余类中的每一个出发，枚举步长取自
\(\{-2,-1,+1,+2\}\)、总位移为零的全部长度 \(2k\) 的字；把式~
\eqref{eq:periodic-operator} 中相应符号相乘并求和。这样得到整数 \(M_k\)，
相减得到 \(F_k\)。当 \(s=1,2,3\) 时，首个正值恰为表中各数，所以引理~
\ref{lem:moment-barrier} 给出 \(R(Q)>8\)。

当 \(s=4\) 时，该字是

\begin{equation}
Q_*=(+,-,-,-,+,-,-,-)
\end{equation}

的一个平移。第~\ref{sec:period-eight-edge} 节证明 \(R(Q_*)=\eta<8\)。我们不
从其前九个超额量的负值作出任何结论。

\subsection{三分律的证明}

合法缺陷数为 \(0,2,4,6,8\)。第~\ref{subsec:all-negative-baseline} 节证明
\(d=0\) 时 \(R=8\)。第~\ref{subsec:four-or-more-defects} 节证明 \(d\geq4\)
时 \(R>8\)。第~\ref{subsec:two-defects} 节证明两缺陷间距为一、二、三时
\(R>8\)，而间距四恰为目标轨道，且 \(R=\eta<8\)。这些情形互不相交且
穷尽所有可能，故定理~\ref{thm:eight-barrier-trichotomy} 得证。

\(Q_*\) 的轨道含四个通量字，对应四个对径点对。因此除平移和反射外，极小
元唯一。全负字是八处唯一的等号情形，因而是唯一的次优轨道。

\subsection{目标相位的手征结构}

目标提升可写为

\begin{equation}
\label{eq:target-antiperiodicity}
\tau=(+,-,+,-,-,+,-,+),       \tau_{i+4}=-\tau_i.
\end{equation}

令 \(T_4(z)\) 为 Bloch 纤维上平移四个点的算子，并令
\(D=\operatorname{diag}((-1)^i)\)。在一般纤维上，

\begin{equation}
(D T_{4}(z))^{2}=zI.
\end{equation}

选择满足 \(\xi^2=z\) 的 \(\xi\)，并定义 \(J_z=\xi^{-1}DT_4(z)\)。把式~
\eqref{eq:target-antiperiodicity} 代入四种跃迁可得

\begin{equation}
\label{eq:chiral-symmetry}
J_z^{2}=I,       J_z H(z)J_z^{-1}=-H(z).
\end{equation}

\(J_z\) 的两个特征空间维数均为四。在手征基下，

\begin{equation}
\label{eq:off-diagonal-square}
H(z)=\begin{pmatrix}0&B\\C&0\end{pmatrix},\qquad H(z)^2=\begin{pmatrix}BC&0\\0&CB\end{pmatrix}.
\end{equation}

在单位圆上 \(C=B^*\)。式~\eqref{eq:off-diagonal-square} 解释了偶特征多项式
和四维平方能带约化。它本身并不证明最优性：合法手征字构成两个二面体轨道，
一个有两个缺陷，一个有六个缺陷；后者由第~\ref{subsec:four-or-more-defects}
节可知位于八之上。
